\documentclass{article}

\usepackage{amsmath,amssymb}
\usepackage{xurl}
\usepackage[hidelinks]{hyperref}
\setlength{\emergencystretch}{2em}

% Shared commands for notes in docs/paper-gaps/.

\DeclareUnicodeCharacter{2080}{\ifmmode{}_{0}\else\textsubscript{0}\fi}
\DeclareUnicodeCharacter{2081}{\ifmmode{}_{1}\else\textsubscript{1}\fi}
\DeclareUnicodeCharacter{2082}{\ifmmode{}_{2}\else\textsubscript{2}\fi}
\DeclareUnicodeCharacter{2099}{\ifmmode{}_{n}\else\textsubscript{n}\fi}

\newcommand{\Lean}{\textsc{Lean}}
\ExplSyntaxOn
\NewDocumentCommand{\leanid}{m}{
  \ifmmode
    \textsf{\detokenize{#1}}
  \else
    \group_begin:
    \urlstyle{sf}
    \tl_set:Nn \l_tmpa_tl {#1}
    \tl_replace_all:Nnn \l_tmpa_tl {\_} {_}
    \exp_args:NV \path \l_tmpa_tl
    \group_end:
  \fi
}
\ExplSyntaxOff
\providecommand{\tr}{\operatorname{tr}}
\newcommand{\tnleanIssueBase}{https://github.com/LionSR/TNLean/issues/}
\newcommand{\tnleanPrBase}{https://github.com/LionSR/TNLean/pull/}
\newcommand{\tnleanDocBase}{https://sirui-lu.com/TNLean/docs/}
\newcommand{\ghissue}[1]{\href{\tnleanIssueBase#1}{GitHub issue \##1}}
\newcommand{\ghpr}[1]{\href{\tnleanPrBase#1}{PR \##1}}
\newcommand{\doclink}[1]{\href{\tnleanDocBase#1.html}{\path{#1}}}

% Dirac notation and the identity operator, as in blueprint/src/macros/common.tex.
% \providecommand keeps note-local definitions authoritative.
\usepackage{dsfont}
\providecommand{\ket}[1]{|#1\rangle}
\providecommand{\bra}[1]{\langle #1|}
\providecommand{\braket}[2]{\langle #1 | #2 \rangle}
\providecommand{\ketbra}[2]{|#1\rangle\!\langle #2|}
\providecommand{\Id}{\mathds{1}}
\providecommand{\one}{\mathds{1}}
\providecommand{\C}{\mathbb{C}}
\providecommand{\MN}[1]{M_{#1}(\C)}
\providecommand{\MD}{M_D(\C)}

% External referencing for paper sources.  The build helper writes sanitized
% auxiliary files under build/paper-aux/ and excludes duplicated source labels.
\usepackage{xr}

% Import a generated paper auxiliary file with a caller-supplied label prefix.
% The candidate paths support builds from docs/paper-gaps/, the repository root,
% and build/paper-aux/ itself.
\newcommand{\importpaperaux}[2]{%
  \IfFileExists{../../build/paper-aux/#2.aux}{%
    \externaldocument[#1]{../../build/paper-aux/#2}%
  }{%
    \IfFileExists{build/paper-aux/#2.aux}{%
      \externaldocument[#1]{build/paper-aux/#2}%
    }{%
      \IfFileExists{#2.aux}{%
        \externaldocument[#1]{#2}%
      }{}%
    }%
  }%
}

% General external reference macros
\newcommand{\paperref}[2]{\ref{#1-#2}}
\newcommand{\papereqref}[2]{(\ref{#1-#2})}

% CPSV16 source (1606.00608 / MPDO-22-12-17-2.tex) external document and shortcuts
\importpaperaux{cpsv16-}{1606.00608/MPDO-22-12-17-2-tnlean}
\newcommand{\cpsvref}[1]{\paperref{cpsv16}{#1}}
\newcommand{\cpsveqref}[1]{\papereqref{cpsv16}{#1}}

% Verdict marker: \gapnote{<kind>}{<status>}. Kind is the policy
% classification (clarification, local-correction, scope-restriction,
% unfaithful, false-source, open-gap); status is open, wip (elimination
% underway), resolved, or historical. Machine-read by the site index;
% prints a verdict line.
\newcommand{\gapnote}[2]{\par\noindent\textbf{Verdict:} #1 (#2).\par}


\title{Closure record for saturation of the area law from MPDO
renormalization fixed points}
\author{}
\date{2026-07-15}

\begin{document}
\maketitle
\gapnote{clarification}{resolved}

\section{Source statement and normalization boundary}

Definition~4.1 of Ref.~\cite{Cirac2016MPDO_arXiv} defines a mixed-state
renormalization fixed point using trace-preserving completely positive maps
$\mathcal S$ and $\mathcal T$. If $M_n(X)$ denotes the $n$-site matrix product
operator with virtual boundary operator $X$, the defining equations are
\begin{align}
    \mathcal S[M_2(X)]
        &=
        M_1(X),
    \label{eq:cpsv16_rfp_sal_data_processing_source_coarse_graining}\\
    \mathcal T[M_1(X)]
        &=
        M_2(X).
    \label{eq:cpsv16_rfp_sal_data_processing_source_fine_graining}
\end{align}
These identities hold for every virtual boundary operator $X$.\footnote{Source
label \texttt{RFPMixedTS},
\path{Papers/1606.00608/MPDO-22-12-17-2.tex:638--660}.}
Proposition \texttt{propsimple} in Appendix~C of
Ref.~\cite{Cirac2016MPDO_arXiv} states that these local renormalization
equations imply source zero correlation length (ZCL) and saturation of the area
law (SAL).\footnote{Source:
\path{Papers/1606.00608/MPDO-22-12-17-2.tex:1331--1341}.}

The entropic convention preceding Definition~4.6 in
Ref.~\cite{Cirac2016MPDO_arXiv} normalizes each finite-chain operator by its
trace.\footnote{Source:
\path{Papers/1606.00608/MPDO-22-12-17-2.tex:792--793}.}
Consequently, Proposition \texttt{propsimple} applies at the following
normalization boundary:
\begin{enumerate}
    \item $\rho^{(N)}(M)$ is positive semidefinite for every $N\geq 1$;
    \item $\tr\rho^{(N)}(M)\ne 0$ for every $N\geq 1$;
    \item $M$ satisfies the local renormalization equations in
    Eqs.~\ref{eq:cpsv16_rfp_sal_data_processing_source_coarse_graining}
    and~\ref{eq:cpsv16_rfp_sal_data_processing_source_fine_graining}.
\end{enumerate}
Positivity does not imply the second condition. The zero tensor generates a
positive semidefinite zero operator at every length, but
\begin{align}
    \tr\rho^{(N)}(0)
        &=
        0.
    \label{eq:cpsv16_rfp_sal_data_processing_zero_tensor_trace}
\end{align}
The nonzero-trace condition is therefore the normalization boundary implicit
in the source's density-operator language.

\section{Area-law calculation}

Let $\sigma^{(N)}_{a+1:b+2}$ denote the normalized density operator on an
$N$-site ring partitioned into consecutive intervals of lengths $a+1$ and
$b+2$. Localizing $\mathcal T$ and $\mathcal S$ to the two intervals gives
channels $\widehat{\mathcal T}_a$ and $\widehat{\mathcal S}_b$ satisfying
\begin{align}
    (\widehat{\mathcal T}_a\otimes\widehat{\mathcal S}_b)
    (\sigma^{(N)}_{a+1:b+2})
        &=
        \sigma^{(N)}_{a+2:b+1}.
    \label{eq:cpsv16_rfp_sal_data_processing_forward_bipartition_channel}
\end{align}
The reverse localized channels satisfy
\begin{align}
    (\widehat{\mathcal S}_a\otimes\widehat{\mathcal T}_b)
    (\sigma^{(N)}_{a+2:b+1})
        &=
        \sigma^{(N)}_{a+1:b+2}.
    \label{eq:cpsv16_rfp_sal_data_processing_reverse_bipartition_channel}
\end{align}
Both are identities of normalized density operators because the corresponding
ring trace is nonzero.

Mutual-information data processing in the two directions gives
\begin{align}
    I(\sigma^{(N)}_{a+2:b+1})
        &\leq
        I(\sigma^{(N)}_{a+1:b+2}),
    \label{eq:cpsv16_rfp_sal_data_processing_forward_mutual_information}\\
    I(\sigma^{(N)}_{a+1:b+2})
        &\leq
        I(\sigma^{(N)}_{a+2:b+1}).
    \label{eq:cpsv16_rfp_sal_data_processing_reverse_mutual_information}
\end{align}
Hence
\begin{align}
    I(\sigma^{(N)}_{a+1:b+2})
        &=
        I(\sigma^{(N)}_{a+2:b+1}).
    \label{eq:cpsv16_rfp_sal_data_processing_mutual_information_equality}
\end{align}
The consecutive-cut identification turns this equality into
\begin{align}
    I_L
        &=
        I_{L+1},
        \qquad
        1\leq L<\left\lfloor\frac{N}{2}\right\rfloor.
    \label{eq:cpsv16_rfp_sal_data_processing_area_law_saturation}
\end{align}

The primary formal declaration is
\leanid{MPOTensor.isSAL_of_isRFPViaTS_of_trace_ne_zero} in
\doclink{TNLean/MPS/MPDO/RFPViaTSSAL}. The channel identities and the
bipartite-to-chain identification are supplied by
\leanid{MPOTensor.tensorMapBoth_bipartitionedNormalizedMPO},
\leanid{MPOTensor.tensorMapBoth_bipartitionedNormalizedMPO_reverse}, and
\leanid{MPOTensor.mutualInfoChain_eq_mutualInformation}.

\section{Source zero correlation length}

Let $\mathcal T_M$ denote the physical-trace transfer operator of $M$. Trace
preservation in the one-to-two renormalization equation gives, for every
virtual matrix $X$,
\begin{align}
    \tr(\mathcal T_M^2X)
        &=
        \tr(M_2(X)),
    \label{eq:cpsv16_rfp_sal_data_processing_transfer_trace_length_two}\\
    \tr(M_2(X))
        &=
        \tr(M_1(X)),
    \label{eq:cpsv16_rfp_sal_data_processing_renormalization_trace_equality}\\
    \tr(M_1(X))
        &=
        \tr(\mathcal T_MX).
    \label{eq:cpsv16_rfp_sal_data_processing_transfer_trace_length_one}
\end{align}
Nondegeneracy of the trace pairing therefore yields
\begin{align}
    \mathcal T_M^2
        &=
        \mathcal T_M.
    \label{eq:cpsv16_rfp_sal_data_processing_physical_trace_idempotence}
\end{align}
If $\mathcal T_M=0$, then
\begin{align}
    \tr\rho^{(1)}(M)
        &=
        \tr(\mathcal T_M)
        =
        0,
    \label{eq:cpsv16_rfp_sal_data_processing_zero_transfer_trace}
\end{align}
contradicting the positive-length nonvanishing hypothesis. Thus
\begin{align}
    \mathcal T_M
        &\ne
        0.
    \label{eq:cpsv16_rfp_sal_data_processing_nonzero_transfer}
\end{align}
Literal physical-trace idempotence follows directly; nonvanishing supplies the
separate scale-invariant relation.

The direct formalization of the literal Definition~4.2 conclusion from
Ref.~\cite{Cirac2016MPDO_arXiv} is
\leanid{MPOTensor.isPhysicalTraceIdempotent_of_isRFPViaTS} in
\doclink{TNLean/MPS/MPDO/RFPViaTS}. Combining it with
\leanid{MPOTensor.isSAL_of_isRFPViaTS_of_trace_ne_zero} gives the source-facing
forward implication: positive semidefinite rings, nonzero trace at every
positive length, and the local renormalization equations of Definition~4.1
imply literal ZCL and SAL. The convenience theorem
\leanid{MPOTensor.isSourceZCL_and_isSAL_of_isRFPViaTS_of_trace_ne_zero}
records the broader scale-invariant ZCL corollary and is not the paper-facing
statement.

\section{Horizontal specialization}

The declaration \leanid{MPOTensor.isSAL_of_isRFPViaTS} gives the horizontal
SAL specialization. Together with
\leanid{MPOTensor.isPhysicalTraceIdempotent_of_isRFPViaTS}, it yields the
paper-facing horizontal corollary. The broader scale-invariant conclusion
follows from \leanid{MPOTensor.isSourceZCL_and_isSAL_of_isRFPViaTS_of_trace_ne_zero}
once the horizontal hypothesis supplies nonvanishing positive-length traces.

The normalized BNT-refined horizontal hypothesis is stronger than the
density-family boundary above. It supplies
\begin{align}
    \tr\rho^{(N)}(M)
        &\ne
        0,
        \qquad N\geq 1,
    \label{eq:cpsv16_rfp_sal_data_processing_horizontal_nonzero_trace}
\end{align}
through
\leanid{MPOTensor.trace_mpo_ne_zero_of_isHorizontalCF_isMPDO_isRFPViaTS}.
The mutual-information argument itself uses no horizontal decomposition.

The literal source canonical-form construction in
Ref.~\cite{Cirac2016MPDO_arXiv} permits an additional all-zero virtual summand
because it allows
\begin{align}
    \sum_k D_k
        &\leq
        D.
    \label{eq:cpsv16_rfp_sal_data_processing_source_virtual_dimension_bound}
\end{align}
\footnote{Source:
\path{Papers/1606.00608/MPDO-22-12-17-2.tex:214--219}.}
This observation does not leave Proposition \texttt{propsimple} open. The
proposition is stated and proved at its density-operator normalization
boundary, while horizontal form provides a stronger sufficient condition for
that boundary.

\section{Classification}

Proposition \texttt{propsimple} of Ref.~\cite{Cirac2016MPDO_arXiv} is complete
under the explicit nonzero-trace density-family condition. Theorem~4.9 of
Ref.~\cite{Cirac2017MPDO_AnnPhys} is partial and partly refuted for independent
reasons. The normal Case-I form of Lemma~C.5 is complete. In Case II, the
conditional projector-controlled channel assembly after common-weight
absorption is complete, but the source's representative-factorization premise
is false under the complete standing hypotheses. Separately, unequal raw
repeated-copy weights refute the literal implication (iv)$\Rightarrow$(v).
The bare source conclusion (ii)$\Rightarrow$(v) is neither proved nor refuted,
because the counterexample to the printed route does not exclude different
blocking channels. The completed forward implication uses no simplicity
hypothesis.

\section{Conclusion}

The source implication from local mixed-state RFP maps to ZCL and SAL is valid
at the normalization boundary implicit in its density-operator convention.
Positivity and local renormalization alone do not exclude the zero family;
nonzero trace at every positive length does. Under that condition, the two
localized channel identities give mutual-information inequalities in opposite
directions, and
Eq.~\ref{eq:cpsv16_rfp_sal_data_processing_mutual_information_equality}
establishes SAL. Trace preservation also gives the literal transfer
idempotence in
Eq.~\ref{eq:cpsv16_rfp_sal_data_processing_physical_trace_idempotence}.
Accordingly, this note records a resolved clarification rather than a change to
the source proposition.

\bibliographystyle{alpha}
\bibliography{references}

\end{document}
